% Noether P09 Korean producer draft — UNCHECKED
% T05-U40; ED0002 lines6836--6844; source 660 LF B / 64111328B2BE86A146300ABAF8AA82329E16D399F569317A6F1CE1E5B9B79A7F
% Pointer v009 / ED0002: B06BE3530D9CF2E82B56FDBA7FE41D5D044DF2425DFA2A059D4939EAA2F7A6C2 / C9A125167ACB33D914EE4374B65AE7CDF0052F568371B8B77B720EA178ABF0E3
% Translation only; example logic remains checker-held.

따라서 영역 $\{\mathfrak G_\eta,\psi\}$는 성질 I과 V를 가지며, $\mathfrak G_\eta$가 그 부분정역이므로 II와 III도 가진다. IV의 성립 여부는 $\psi$의 선택에 달려 있다. 예를 들어 $\psi=\frac{1}{n\cdot\eta}$이면 $\frac1n$이 영역에 속하므로 IV가 성립하지 않는다. 반면 § 2에 따라 $\psi=\frac1\eta$일 때나 $\psi=\frac{\eta}{n}$일 때에는 IV가 성립한다. 실제로 후자의 경우 $\{\mathfrak G_\eta,\psi\}$의 어떤 함수가 대수적 수를 나타내면 $\eta=0$으로 놓아도 그 값이 유지된다. 그러면 그 함수는 $\mathfrak G_\eta$의 함수로 바뀌고, 따라서 대수적 정수로 바뀐다.
